\section{Elementary Koenigs functions at \texorpdfstring{$r=2$ and $r=4$}{r=2 and r=4}}
\label{a:app:integrable}

These cases lie outside subcritical branching, where $r=2p<1$, but complete the
classical picture of integrable logistic linearisations.  They also locate the two
elementary examples within the classification used in the main text.

\subsection{Case \texorpdfstring{$r=2$}{r=2} (logarithmic)}

For $r=2$, the map is $f_2(x)=2x(1-x)$ and the fixed point $0$ has the
repelling multiplier $2$.  Its linearisation is therefore of Poincar\'e type rather
than attracting-Koenigs type, but it still has an elementary form.

The affine change $u=1-2x$ conjugates the dynamics to $u\mapsto u^2$.
Linearising this monomial by the logarithm and returning to $x$ gives
\begin{equation}
\label{a:eq:phi2}
\psi_2(x)=-\frac12\ln(1-2x),
\end{equation}
normalised so that $\psi_2(0)=0$ and $\psi_2'(0)=1$.

\textbf{Verification.}
\begin{align*}
\psi_2\bigl(f_2(x)\bigr)
&=
-\frac12\ln\bigl(1-2\cdot 2x(1-x)\bigr)
=
-\frac12\ln\bigl((1-2x)^2\bigr)
=
-\ln(1-2x)
=
2\psi_2(x).
\end{align*}

\subsection{Case \texorpdfstring{$r=4$}{r=4} (inverse trigonometric)}

For $r=4$, the fixed point $0$ of $f_4(x)=4x(1-x)$ has the repelling multiplier
$4$, so the linearisation is again of Poincar\'e type.  Writing
$x=\sin^2\theta$ gives $f_4(x)=\sin^2(2\theta)$ and therefore doubles the angle.
This suggests
\begin{equation}
\label{a:eq:phi4}
\psi_4(x)=\bigl(\arcsin\sqrt{x}\bigr)^2
\end{equation}
satisfies $\psi_4(f_4(x))=4\psi_4(x)$ on a suitable real domain, and $\psi_4(x)\sim x$ as $x\to0$, so $\psi_4'(0)=1$.

\textbf{Verification.}
Using $\arcsin(2u\sqrt{1-u^2})=2\arcsin u$ with $u=\sqrt{x}$ (on an appropriate interval),
\begin{align*}
\psi_4\bigl(f_4(x)\bigr)
&=
\bigl(\arcsin\sqrt{4x(1-x)}\bigr)^2
=
\bigl(2\arcsin\sqrt{x}\bigr)^2
=
4\psi_4(x).
\end{align*}

\subsection{Affine conjugacy to \texorpdfstring{$z^2+c$}{z squared plus c}}

\begin{lemma}
\label{a:lem:affine}
The logistic map $f_r(x)=rx(1-x)$ is affinely conjugate to $P_c(z)=z^2+c$ with $c=(2r-r^2)/4$.
\end{lemma}

\begin{proof}[Proof sketch]
Seek $h(x)=ax+b$ with $h\circ f_r\circ h^{-1}=P_c$. The substitution $x=-z/r+\tfrac12$ (for $r\neq0$) yields a monic quadratic; matching coefficients gives $c=(2r-r^2)/4$.
\end{proof}

Under this conjugacy, $r=2$ corresponds to the monomial parameter $c=0$, while
$r=4$ corresponds to the Chebyshev parameter $c=-2$; the latter is also reached
at $r=-2$.  Neither value lies in the image $c\in(0,\tfrac14)$ of $r\in(0,1)$.
As \Cref{a:rem:empty} explains, this parameter count is a consistency check rather
than the proof.  The decisive fact is structural: both elementary cases sit over
repelling fixed points, which the attracting window excludes.

\subsection{The worked examples inside the Becker--Bergweiler classification}

The classification in \Cref{a:thm:BB-main} places every differentially algebraic
Schr\"oder function over a repelling fixed point.  Such a function is a M\"obius
transformation of exponential, cosine or Weierstrass-$\wp$ material, while the
underlying map is M\"obius-conjugate to $z^{\pm d}$, $\pm T_d$, or a Latt\`es map.
The two derivations above realise the exponential and cosine families explicitly.

\emph{Case $r=2$.} Inverting \eqref{a:eq:phi2}: $t=-\tfrac12\ln(1-2x)$ gives
\[
\varphi_2(t)=\psi_2^{-1}(t)=\frac{1-e^{-2t}}{2},
\]
an affine, and hence M\"obius, image of $\exp(-2t)$: the family
$\exp(\alpha t^{\rho})$ with $\rho=1$ and $\alpha=-2$.  Correspondingly,
$u=1-2x$ conjugates $f_2$ to $u\mapsto u^2$, and the fixed point has the
repelling multiplier $2$, as \Cref{a:thm:BB-main}(i) requires.

\emph{Case $r=4$.} Inverting \eqref{a:eq:phi4}: $t=(\arcsin\sqrt{x})^2$ gives
\[
\varphi_4(t)=\psi_4^{-1}(t)=\sin^2\!\sqrt{t}=\frac{1-\cos\bigl(2\sqrt{t}\bigr)}{2},
\]
a M\"obius image of $\cos(\alpha t^{\rho}+\beta)$ with $\rho=\tfrac12$,
$\alpha=2$ and $\beta=0$.  The fractional exponent is admissible because
$\cos(2\sqrt{t})$ is entire in $t$.  Correspondingly, $u=1-2x$ conjugates
$f_4$ to the second Chebyshev polynomial $T_2(u)=2u^2-1$, and the fixed point has
the repelling multiplier $4$.

Both examples therefore occupy the positions prescribed by
\Cref{a:thm:BB-main}: repelling fixed points in the exponential and cosine
families.  For $r\in(0,1)$ the fixed point is attracting, no classified elementary
case is available, and \Cref{a:thm:main-ht} gives hypertranscendence of $\psi_r$.
That distinction underpins the main-text claim for the basin value
$A(p)=2\psi_{2p}(\tfrac12)$.
